\documentclass[10pt,a4paper]{amsart}
\usepackage{amsmath,amssymb,amsthm,graphicx}
% Retained minimal preamble for the Stacks Project chapter flat.tex (More on Flatness).
% Verbatim excerpt of the chapter preamble: Xy-pic setup (native preamble lines 14-18) and the native
% theorem declarations with the mathematical macros (native preamble lines 154-217). Copied mechanically;
% nothing added, removed or redefined. The wrapper class and the amsmath/amssymb/amsthm/graphicx base
% packages are supplied by the assembler, so the class/externaldocument part of the native preamble is omitted.
% For commutative diagrams we use Xy-pic
\usepackage[all]{xy}

% We use 2cell for 2-commutative diagrams.
\xyoption{2cell}
% Theorem environments.
%
\theoremstyle{plain}
\newtheorem{theorem}[subsection]{Theorem}
\newtheorem{proposition}[subsection]{Proposition}
\newtheorem{lemma}[subsection]{Lemma}

\theoremstyle{definition}
\newtheorem{definition}[subsection]{Definition}
\newtheorem{example}[subsection]{Example}
\newtheorem{exercise}[subsection]{Exercise}
\newtheorem{situation}[subsection]{Situation}

\theoremstyle{remark}
\newtheorem{remark}[subsection]{Remark}
\newtheorem{remarks}[subsection]{Remarks}

\numberwithin{equation}{subsection}

% Macros
%
\def\lim{\mathop{\mathrm{lim}}\nolimits}
\def\colim{\mathop{\mathrm{colim}}\nolimits}
\def\Spec{\mathop{\mathrm{Spec}}}
\def\Hom{\mathop{\mathrm{Hom}}\nolimits}
\def\Ext{\mathop{\mathrm{Ext}}\nolimits}
\def\SheafHom{\mathop{\mathcal{H}\!\mathit{om}}\nolimits}
\def\SheafExt{\mathop{\mathcal{E}\!\mathit{xt}}\nolimits}
\def\Sch{\mathit{Sch}}
\def\Mor{\mathop{\mathrm{Mor}}\nolimits}
\def\Ob{\mathop{\mathrm{Ob}}\nolimits}
\def\Sh{\mathop{\mathit{Sh}}\nolimits}
\def\NL{\mathop{N\!L}\nolimits}
\def\CH{\mathop{\mathrm{CH}}\nolimits}
\def\proetale{{pro\text{-}\acute{e}tale}}
\def\etale{{\acute{e}tale}}
\def\QCoh{\mathit{QCoh}}
\def\Ker{\mathop{\mathrm{Ker}}}
\def\Im{\mathop{\mathrm{Im}}}
\def\Coker{\mathop{\mathrm{Coker}}}
\def\Coim{\mathop{\mathrm{Coim}}}

% Boxtimes
%
\DeclareMathSymbol{\boxtimes}{\mathbin}{AMSa}{"02}

%
% Macros for moduli stacks/spaces
%
\def\QCohstack{\mathcal{QC}\!\mathit{oh}}
\def\Cohstack{\mathcal{C}\!\mathit{oh}}
\def\Spacesstack{\mathcal{S}\!\mathit{paces}}
\def\Quotfunctor{\mathrm{Quot}}
\def\Hilbfunctor{\mathrm{Hilb}}
\def\Curvesstack{\mathcal{C}\!\mathit{urves}}
\def\Polarizedstack{\mathcal{P}\!\mathit{olarized}}
\def\Complexesstack{\mathcal{C}\!\mathit{omplexes}}
% \Pic is the operator that assigns to X its picard group, usage \Pic(X)
% \Picardstack_{X/B} denotes the Picard stack of X over B
% \Picardfunctor_{X/B} denotes the Picard functor of X over B
\def\Pic{\mathop{\mathrm{Pic}}\nolimits}
\def\Picardstack{\mathcal{P}\!\mathit{ic}}
\def\Picardfunctor{\mathrm{Pic}}
\def\Deformationcategory{\mathcal{D}\!\mathit{ef}}

% --- appended 2026-10-12 (Chief): the project own definitions of the reference/slogan/history
% display environments, copied mechanically from the shared Stacks preamble lines 8-12
% (\usepackage{verbatim} plus the three \newenvironment lines, which the chapter preamble keeps
% verbatim from the project preamble). Nothing added, removed or redefined: these environments
% carry no mathematical content and the project itself typesets them as comments. Without these
% three lines a counted statement whose own text uses them fails with Environment slogan undefined.
\usepackage{verbatim}
\newenvironment{reference}{\comment}{\endcomment}
\newenvironment{slogan}{\comment}{\endcomment}
\newenvironment{history}{\comment}{\endcomment}
\usepackage{cleveref}
\crefname{originallemma}{original lemma}{original lemmas}
\Crefname{originallemma}{Original lemma}{Original lemmas}
\makeatletter
\newlabel{divisors-lemma-flat-base-change-blowing-up}{{0805}{source}{}{sourcepointer-1}{}}
\newlabel{divisors-lemma-flat-base-change-blowing-up@cref}{{[originallemma][0][]0805}{[1][0][]source}{}{}{}}
\makeatother
\crefname{originallemma}{original lemma}{original lemmas}
\Crefname{originallemma}{Original lemma}{Original lemmas}
\makeatletter
\newlabel{divisors-lemma-blowing-up-two-ideals}{{080A}{source}{}{sourcepointer-2}{}}
\newlabel{divisors-lemma-blowing-up-two-ideals@cref}{{[originallemma][0][]080A}{[1][0][]source}{}{}{}}
\makeatother
\crefname{originallemma}{original lemma}{original lemmas}
\Crefname{originallemma}{Original lemma}{Original lemmas}
\makeatletter
\newlabel{divisors-lemma-strict-transform-flat}{{080F}{source}{}{sourcepointer-3}{}}
\newlabel{divisors-lemma-strict-transform-flat@cref}{{[originallemma][0][]080F}{[1][0][]source}{}{}{}}
\makeatother
\crefname{originallemma}{original lemma}{original lemmas}
\Crefname{originallemma}{Original lemma}{Original lemmas}
\makeatletter
\newlabel{divisors-lemma-strict-transform-composition-blowups}{{080I}{source}{}{sourcepointer-4}{}}
\newlabel{divisors-lemma-strict-transform-composition-blowups@cref}{{[originallemma][0][]080I}{[1][0][]source}{}{}{}}
\makeatother
\crefname{originallemma}{original lemma}{original lemmas}
\Crefname{originallemma}{Original lemma}{Original lemmas}
\makeatletter
\newlabel{morphisms-lemma-base-change-module-flat}{{01U8}{source}{}{sourcepointer-5}{}}
\newlabel{morphisms-lemma-base-change-module-flat@cref}{{[originallemma][0][]01U8}{[1][0][]source}{}{}{}}
\makeatother
\crefname{originalreference}{original reference}{original references}
\Crefname{originalreference}{Original reference}{Original references}
\makeatletter
\newlabel{remark-successive-blowups}{{080Y}{source}{}{sourcepointer-6}{}}
\newlabel{remark-successive-blowups@cref}{{[originalreference][0][]080Y}{[1][0][]source}{}{}{}}
\makeatother
\crefname{originallemma}{original lemma}{original lemmas}
\Crefname{originallemma}{Original lemma}{Original lemmas}
\makeatletter
\newlabel{divisors-lemma-characterize-effective-Cartier-divisor}{{01WS}{source}{}{sourcepointer-7}{}}
\newlabel{divisors-lemma-characterize-effective-Cartier-divisor@cref}{{[originallemma][0][]01WS}{[1][0][]source}{}{}{}}
\makeatother
\crefname{originallemma}{original lemma}{original lemmas}
\Crefname{originallemma}{Original lemma}{Original lemmas}
\makeatletter
\newlabel{lemma-push-ideal}{{0813}{source}{}{sourcepointer-8}{}}
\newlabel{lemma-push-ideal@cref}{{[originallemma][0][]0813}{[1][0][]source}{}{}{}}
\makeatother
\crefname{originallemma}{original lemma}{original lemmas}
\Crefname{originallemma}{Original lemma}{Original lemmas}
\makeatletter
\newlabel{lemma-etale-flat-up-down}{{05B9}{source}{}{sourcepointer-9}{}}
\newlabel{lemma-etale-flat-up-down@cref}{{[originallemma][0][]05B9}{[1][0][]source}{}{}{}}
\makeatother
\crefname{originallemma}{original lemma}{original lemmas}
\Crefname{originallemma}{Original lemma}{Original lemmas}
\makeatletter
\newlabel{morphisms-lemma-flat-permanence}{{02JZ}{source}{}{sourcepointer-10}{}}
\newlabel{morphisms-lemma-flat-permanence@cref}{{[originallemma][0][]02JZ}{[1][0][]source}{}{}{}}
\makeatother
\title{Etale local flatness after admissible blowups globalises to one admissible blowup making the strict transform flat}
\author{The Stacks Project Authors}
\date{stacks-project:flat.tex:local-snapshot-20261011; GNU FDL 1.2 or later}
\begin{document}
\maketitle
\noindent\textit{Editorial note.} All delivered mathematical text is the author's own, copied from the pinned Stacks Project chapter file flat.tex; no statement or proof step was rewritten, reordered or removed.\par\medskip
\noindent\textit{Editorial note.} The preamble is the chapter's own retained minimal preamble (sources/stacks-flat/retained-preamble.tex), a verbatim excerpt of the native Xy-pic setup and of the native theorem declarations and mathematical macros, followed by the project own definitions of the reference, slogan and history display environments (verbatim package plus three newenvironment lines) that the retained file carries at its end and that the project itself typesets as comments. No macro was invented or redefined, and the wrapper class and the amsmath/amssymb/amsthm/graphicx base packages are supplied by the assembler.\par\medskip
\noindent\textit{Editorial note.} The counted statement and proof contain no cross reference, no citation and no included file, so the delivered body is self contained.\par\medskip
\noindent\textit{Editorial note.} Canonical location of the counted proposition: Stacks tag 0814 (sources/stacks-flat/pinned-original-source.tex line 8904).\par\medskip
\noindent\textit{Allowed standard prerequisites (editorial).}\begin{itemize}
\item Strict transforms and flatness under blowups and blowup composition. Divisors, Lemmas labelled divisors-lemma-strict-transform-flat and divisors-lemma-strict-transform-composition-blowups, cited at proof lines 8958-8961; Divisors, Lemma labelled divisors-lemma-flat-base-change-blowing-up, cited at proof lines 8942-8944. Tracks the strict transform of the module through the base change and the additional blowup.
\item Comparison of ideals in two blowups. Divisors, Lemma labelled divisors-lemma-blowing-up-two-ideals, cited at proof lines 8946-8948 and 9002-9004. Produces the commutative square of blowups and shows the base change of the final blowup factors through the prescribed one.
\item Successive blowups and reduction to the affine divisor case. Same chapter, Remark labelled remark-successive-blowups, cited at proof lines 8974-8978; Divisors, Lemma labelled divisors-lemma-characterize-effective-Cartier-divisor, cited at proof lines 8978-8980. Reduces to S = Spec(A) with S minus U the zero locus of a nonzerodivisor.
\item Splitting an ideal after an etale base change. Same chapter, Lemma labelled lemma-push-ideal, cited at proof lines 8984-8986. Produces the finite type ideals I\_i on S whose pullbacks split off the given ideals J\_i.
\item Flatness of a module and its pullback along an etale cover. Same chapter, Lemma labelled lemma-etale-flat-up-down, cited at proof lines 9008-9010; Morphisms, Lemma labelled morphisms-lemma-flat-permanence, cited at proof lines 9010-9012; Morphisms, Lemma labelled morphisms-lemma-base-change-module-flat, cited at proof lines 8968-8970. Descends flatness from the etale cover S\_i back to S.
\end{itemize}\medskip
\ifdefined\hypertarget\hypertarget{sourcepointer-1}{}\fi
\noindent\textit{Source reference.} Original lemma 0805: sources/stacks-chow/tags.txt, source lines 10170--10170. Cross-chapter standard prerequisite located in the project tag list staged in the local snapshot; invoked at proof line 8940 of the counted proof. It is retained only as a navigation pointer; every load-bearing step is inside the delivered span.\par\medskip
\ifdefined\hypertarget\hypertarget{sourcepointer-2}{}\fi
\noindent\textit{Source reference.} Original lemma 080A: sources/stacks-chow/tags.txt, source lines 10175--10175. Cross-chapter standard prerequisite located in the project tag list staged in the local snapshot; invoked at proof line 8944 of the counted proof. It is retained only as a navigation pointer; every load-bearing step is inside the delivered span.\par\medskip
\ifdefined\hypertarget\hypertarget{sourcepointer-3}{}\fi
\noindent\textit{Source reference.} Original lemma 080F: sources/stacks-chow/tags.txt, source lines 10180--10180. Cross-chapter standard prerequisite located in the project tag list staged in the local snapshot; invoked at proof line 8953 of the counted proof. It is retained only as a navigation pointer; every load-bearing step is inside the delivered span.\par\medskip
\ifdefined\hypertarget\hypertarget{sourcepointer-4}{}\fi
\noindent\textit{Source reference.} Original lemma 080I: sources/stacks-chow/tags.txt, source lines 10183--10183. Cross-chapter standard prerequisite located in the project tag list staged in the local snapshot; invoked at proof line 8954 of the counted proof. It is retained only as a navigation pointer; every load-bearing step is inside the delivered span.\par\medskip
\ifdefined\hypertarget\hypertarget{sourcepointer-5}{}\fi
\noindent\textit{Source reference.} Original lemma 01U8: sources/stacks-chow/tags.txt, source lines 2401--2401. Cross-chapter standard prerequisite located in the project tag list staged in the local snapshot; invoked at proof line 8970 of the counted proof. It is retained only as a navigation pointer; every load-bearing step is inside the delivered span.\par\medskip
\ifdefined\hypertarget\hypertarget{sourcepointer-6}{}\fi
\noindent\textit{Source reference.} Original reference 080Y: sources/stacks-flat/pinned-original-source.tex, source lines 8611--8611. Same-chapter pointer into the author environment remark at source line 8611; invoked at proof line 8972 of the counted proof. It is retained only as a navigation pointer; every load-bearing step is inside the delivered span.\par\medskip
\ifdefined\hypertarget\hypertarget{sourcepointer-7}{}\fi
\noindent\textit{Source reference.} Original lemma 01WS: sources/stacks-chow/tags.txt, source lines 2493--2493. Cross-chapter standard prerequisite located in the project tag list staged in the local snapshot; invoked at proof line 8979 of the counted proof. It is retained only as a navigation pointer; every load-bearing step is inside the delivered span.\par\medskip
\ifdefined\hypertarget\hypertarget{sourcepointer-8}{}\fi
\noindent\textit{Source reference.} Original lemma 0813: sources/stacks-flat/pinned-original-source.tex, source lines 8876--8876. Same-chapter pointer into the author environment lemma at source line 8876; invoked at proof line 8984 of the counted proof. It is retained only as a navigation pointer; every load-bearing step is inside the delivered span.\par\medskip
\ifdefined\hypertarget\hypertarget{sourcepointer-9}{}\fi
\noindent\textit{Source reference.} Original lemma 05B9: sources/stacks-flat/pinned-original-source.tex, source lines 124--124. Same-chapter pointer into the author environment lemma at source line 124; invoked at proof line 9011 of the counted proof. It is retained only as a navigation pointer; every load-bearing step is inside the delivered span.\par\medskip
\ifdefined\hypertarget\hypertarget{sourcepointer-10}{}\fi
\noindent\textit{Source reference.} Original lemma 02JZ: sources/stacks-chow/tags.txt, source lines 3354--3354. Cross-chapter standard prerequisite located in the project tag list staged in the local snapshot; invoked at proof line 9012 of the counted proof. It is retained only as a navigation pointer; every load-bearing step is inside the delivered span.\par\medskip
\begin{lemma}
\label{lemma-flatten-module-etale-localize}
Let $S$ be a quasi-compact and quasi-separated scheme.
Let $X \to S$ be a morphism of schemes.
Let $\mathcal{F}$ be a quasi-coherent module on $X$.
Let $U \subset S$ be a quasi-compact open.
Assume there exist finitely many commutative diagrams
$$
\xymatrix{
& X_i \ar[r]_{j_i} \ar[d] & X \ar[d] \\
S_i^* \ar[r] & S_i \ar[r]^{e_i} & S
}
$$
where
\begin{enumerate}
\item $e_i : S_i \to S$ are quasi-compact \'etale morphisms and
$S = \bigcup e_i(S_i)$,
\item $j_i : X_i \to X$ are \'etale morphisms and
$X = \bigcup j_i(X_i)$,
\item $S^*_i \to S_i$ is an $e_i^{-1}(U)$-admissible blowup
such that the strict transform $\mathcal{F}_i^*$ of $j_i^*\mathcal{F}$
is flat over $S^*_i$.
\end{enumerate}
Then there exists a $U$-admissible blowup $S' \to S$ such that
the strict transform of $\mathcal{F}$ is flat over $S'$.
\end{lemma}

\begin{proof}
We claim that the hypotheses of the lemma are preserved under $U$-admissible
blowups. Namely, suppose $b : S' \to S$ is a $U$-admissible blowup in
the quasi-coherent sheaf of ideals $\mathcal{I}$. Moreover, let $S^*_i \to S_i$
be the blowup in the quasi-coherent sheaf of ideals $\mathcal{J}_i$.
Then the collection of morphisms $e'_i : S'_i = S_i \times_S S' \to S'$
and $j'_i : X_i' = X_i \times_S S' \to X \times_S S'$ satisfy conditions
(1), (2), (3) for the strict transform $\mathcal{F}'$ of $\mathcal{F}$
relative to the blowup $S' \to S$. First, observe that $S_i'$ is the
blowup of $S_i$ in the pullback of $\mathcal{I}$, see
Divisors, Lemma \ref{divisors-lemma-flat-base-change-blowing-up}.
Second, consider the blowup
$S_i^{\prime *} \to S_i'$ of $S_i'$ in the pullback of the ideal
$\mathcal{J}_i$.
By Divisors, Lemma \ref{divisors-lemma-blowing-up-two-ideals}
we get a commutative diagram
$$
\xymatrix{
S_i^{\prime *} \ar[r] \ar[rd] \ar[d] & S'_i \ar[d] \\
S_i^* \ar[r] & S_i
}
$$
and all the morphisms in the diagram above are blowups. Hence by
Divisors, Lemmas \ref{divisors-lemma-strict-transform-flat} and
\ref{divisors-lemma-strict-transform-composition-blowups}
we see
\begin{align*}
& \text{ the strict transform of }(j'_i)^*\mathcal{F}'\text{ under }
S_i^{\prime *} \to S_i' \\
= &
\text{ the strict transform of }j_i^*\mathcal{F}\text{ under }
S_i^{\prime *} \to S_i \\
= &
\text{ the strict transform of }\mathcal{F}_i'\text{ under }
S_i^{\prime *} \to S_i' \\
= &
\text{ the pullback of }\mathcal{F}_i^*\text{ via }
X_i \times_{S_i} S_i^{\prime *} \to X_i
\end{align*}
which is therefore flat over $S_i^{\prime *}$
(Morphisms, Lemma \ref{morphisms-lemma-base-change-module-flat}).
Having said this, we see that all observations of
Remark \ref{remark-successive-blowups} apply to the
problem of finding a $U$-admissible blowup such that the
strict transform of $\mathcal{F}$ becomes flat over the base
under assumptions as in the lemma. In particular, we may assume
that $S \setminus U$ is the support of an effective Cartier divisor
$D \subset S$. Another application of Remark \ref{remark-successive-blowups}
combined with
Divisors, Lemma \ref{divisors-lemma-characterize-effective-Cartier-divisor}
shows we may assume that $S = \Spec(A)$ and $D = \Spec(A/(a))$
for some nonzerodivisor $a \in A$.

\medskip\noindent
Pick an $i$ and $s \in S_i$. Lemma \ref{lemma-push-ideal}
implies we can find an open neighbourhood $s \in W_i \subset S_i$
and a finite type quasi-coherent ideal $\mathcal{I} \subset \mathcal{O}_S$
such that $\mathcal{I} \cdot \mathcal{O}_{W_i} = \mathcal{J}_i \mathcal{J}'_i$
for some finite type quasi-coherent ideal
$\mathcal{J}'_i \subset \mathcal{O}_{W_i}$
and such that $V(\mathcal{I}) \subset V(a) = S \setminus U$.
Since $S_i$ is quasi-compact we can replace $S_i$ by a finite collection
$W_1, \ldots, W_n$ of these opens and assume that for each $i$ there exists
a quasi-coherent sheaf of ideals $\mathcal{I}_i \subset \mathcal{O}_S$ such
that $\mathcal{I}_i \cdot \mathcal{O}_{S_i} = \mathcal{J}_i \mathcal{J}'_i$
for some finite type quasi-coherent ideal
$\mathcal{J}'_i \subset \mathcal{O}_{S_i}$.
As in the discussion of the first paragraph of the proof, consider the
blowup $S'$ of $S$ in the product $\mathcal{I}_1 \ldots \mathcal{I}_n$
(this blowup is $U$-admissible by construction). The base change of $S' \to S$
to $S_i$ is the blowup in
$$
\mathcal{J}_i \cdot
\mathcal{J}'_i \mathcal{I}_1 \ldots \hat{\mathcal{I}_i} \ldots \mathcal{I}_n
$$
which factors through the given blowup $S_i^* \to S_i$
(Divisors, Lemma \ref{divisors-lemma-blowing-up-two-ideals}). In the notation
of the diagram above this means that $S_i^{\prime *} = S_i'$. Hence
after replacing $S$ by $S'$ we arrive in the situation that
$j_i^*\mathcal{F}$ is flat over $S_i$. Hence $j_i^*\mathcal{F}$ is flat
over $S$, see
Lemma \ref{lemma-etale-flat-up-down}.
By Morphisms, Lemma \ref{morphisms-lemma-flat-permanence}
we see that $\mathcal{F}$ is flat over $S$.
\end{proof}

\end{document}
